%% =====================================================================
%%  Unit 12 — I can swim!  (Dire ce qu'on sait faire)
%%  Pistes : 12.01–12.04 dialogues · 12.05–12.06 vocabulaire · 12.07–12.08 grammaire
%%           12.09 prononciation · 12.10–12.13 exercices (textes en annexe)
%% =====================================================================

\newcommand{\otIllusUnXII}{\otIllusDuo{swim/104/86/15/0.10,ski/190/92/20/0.09,sing/122/14/13/0.08,
  dance/150/8/11/0.08,run/196/22/12/0.07,cook/100/50/12/0.07}{emma}{Can you swim?}{lea}{Yes, I can!}}

% Tableau des talents (exercices 2 et 3)
\newcommand{\otTalents}{%
  \begin{center}\otTableau
  \begin{tabular}{@{}l@{\hspace{10mm}}c@{\hspace{10mm}}c@{\hspace{10mm}}c@{\hspace{10mm}}c@{}}
   & \otTH{swim} & \otTH{ski} & \otTH{cook} & \otTH{play the guitar}\\\midrule
  \otNomDit{lea}  & \otCheck & \otCheck & \otCheck & \otCroix\\[1mm]
  \otNomDit{jack} & \otCheck & \otCroix & \otCroix & \otCheck\\[1mm]
  \otNomDit{emma} & \otCroix & \otCheck & \otCheck & \otCroix\\\bottomrule
  \end{tabular}\end{center}}

\unite{12}{I can swim!}{Dire ce qu'on sait faire}{otIllusUnXII}

\situation{Épisode 12}{Samedi, journée sportive au parc pour toute l'équipe de \en{Bright Ideas} :
le coach veut savoir qui sait faire quoi. Le soir, Sam a un problème : la chanteuse de son groupe
est malade, et le concert a lieu vendredi…}

\begin{objectifs}
\item dire ce que vous savez faire et ce que vous ne savez pas faire ;
\item poser des questions avec \en{can} et y répondre : \en{Yes, I can. / No, I can't.} ;
\item préciser votre niveau : \en{well}, \en{very well}, \en{not very well}, \en{not at all} ;
\item parler de sports et de loisirs ;
\item distinguer à l'oral \en{can} \phon{kən} et \en{can't} \phon{kænt}.
\end{objectifs}

\auprogramme
  {Parler de ses capacités et de ses loisirs}
  {\en{can}, \en{can't}\par\vspace{0.6mm}Les adverbes : \en{well}, \en{very well}}
  {Sports, loisirs, talents : 29 expressions}
  {\en{can} ou \en{can't} ?}
  {Le sport aux États-Unis}

\semaine{casque/otNavy/Dialogues/Écouter et lire les deux scènes,
  bulle/otGreenDark/Vocabulaire/Apprendre les 29 expressions,
  blocs/otViolet/Grammaire/can et les adverbes,
  sons/otOrangeDark/Sons et culture/can ou can't et les sports,
  crayon/otNavy/Exercices/12 exercices corrigés,
  coche/otGreenDark/Bilan/L'essentiel et le test express}

\clearpage
%% ---------------------------------------------------------------------
\rubrique{ecouter}{Des talents cachés}

\begin{dialogue}{1}{Can you play basketball?}{Samedi, 10 heures, dans un parc de Brooklyn. Un coach
  organise la journée sportive du studio.}{12.01}{12.02}
\dit{coach}{OK, team! \hl{Can you play basketball}?\note{1}}{Bon, l'équipe ! Vous savez jouer au basket ?}
\dit{jack}{\hl{Yes, I can}! I play every Sunday.}{Oui ! Je joue tous les dimanches.}
\dit{emma}{I \hl{can't play} basketball, but I \hl{can run} very fast!\note{2}}{Je ne sais pas jouer au basket, mais je cours très vite !}
\dit{coach}{Great! And you, Léa? \hl{Can you swim}?}{Super ! Et vous, Léa ? Vous savez nager ?}
\dit{lea}{Yes, I can. I swim \hl{very well}.\note{3}}{Oui. Je nage très bien.}
\dit{coach}{\hl{Can you ride a bike}?}{Vous savez faire du vélo ?}
\dit{lea}{Of course! In Lyon, I ride my bike every day.}{Bien sûr ! À Lyon, je fais du vélo tous les jours.}
\dit{jack}{Ha! Mr.\ Carter \hl{can't swim}!}{Ha ! Mr.\ Carter ne sait pas nager !}
\dit{carter}{That's not true, Jack! I can swim\ldots{} but \hl{not very well}!}{C'est faux, Jack ! Je sais nager… mais pas très bien !}
\end{dialogue}

\begin{comprendre}
\pointnote{1} \en{\textbf{Can you play basketball?}} \en{can} exprime la capacité : « savoir »,
  « pouvoir ». Pour une question, on le place devant le sujet ; on répond \en{Yes, I can} ou
  \en{No, I can't}.
\pointnote{2} \en{\textbf{I can't play…, but I can run…}} \en{can't} (= \en{cannot}) est la forme
  négative. \en{can} est suivi de la base verbale, \textbf{sans} \en{to} : \en{I can run}.
\pointnote{3} \en{\textbf{I swim very well.}} \en{well} (« bien ») est un adverbe : il se place après
  le verbe (et son complément). \en{very well} = très bien ; \en{not very well} = pas très bien.
\end{comprendre}

\begin{dialogue}{2}{Can you sing?}{Samedi soir, chez Sam. Maya est là. Sam a l'air inquiet.}{12.03}{12.04}
\dit{sam}{Léa, \hl{can you sing}? Our singer is sick.}{Léa, tu sais chanter ? Notre chanteuse est malade.}
\dit{lea}{Me? I \hl{can sing}, but \hl{not very well}!}{Moi ? Je sais chanter, mais pas très bien !}
\dit{maya}{\hl{Can you play the piano}?\note{4}}{Tu sais jouer du piano ?}
\dit{lea}{Yes, I can. I play \hl{quite well}.}{Oui. Je joue plutôt bien.}
\dit{sam}{Perfect! \hl{Can you come} on Friday?\note{5}}{Parfait ! Tu peux venir vendredi ?}
\dit{lea}{Yes, I can. What time?}{Oui. À quelle heure ?}
\dit{sam}{At eight. And can you bring some snacks?}{À huit heures. Et tu peux apporter de quoi grignoter ?}
\dit{lea}{Ha ha! \hl{Of course I can}!}{Ha ha ! Bien sûr !}
\end{dialogue}

\begin{comprendre}
\pointnote{4} \en{\textbf{play the piano}} On dit \en{play the piano}, \en{play the guitar} (avec
  \en{the}) pour un instrument, mais \en{play tennis}, \en{play basketball} (sans article) pour un sport.
\pointnote{5} \en{\textbf{Can you come?}} \en{can} sert aussi à demander si quelqu'un est disponible
  (« tu peux venir ? ») ou à demander un service (\en{Can you bring…?}).
\end{comprendre}

%% ---------------------------------------------------------------------
\rubrique{mots}{Les mots de l'unité}

Voici les 29 expressions de l'unité. Retenez les verbes avec leur complément : on \en{plays}
un sport ou un instrument, on \en{rides} un vélo, on \en{drives} une voiture.

\begin{vocabulaire}{Sports et loisirs}{12.05}
\categorie{Les sports}
\mot{swim}{swɪm}{nager}
\mot{run}{rʌn}{courir}
\mot{ride a bike}{ˌraɪd ə ˈbaɪk}{faire du vélo}
\mot{ski}{skiː}{skier}
\mot{play tennis}{ˌpleɪ ˈtenɪs}{jouer au tennis}
\mot{play basketball}{ˌpleɪ ˈbæskɪtbɑːl}{jouer au basket}
\mot{play soccer}{ˌpleɪ ˈsɑːkɚ}{jouer au football \registre{américain}}
\categorie{Les autres talents}
\mot{sing}{sɪŋ}{chanter}
\mot{dance}{dæns}{danser}
\mot{play the piano}{ˌpleɪ ðə piˈænoʊ}{jouer du piano}
\mot{play the guitar}{ˌpleɪ ðə ɡɪˈtɑːr}{jouer de la guitare}
\mot{draw}{drɑː}{dessiner}
\mot{drive a car}{ˌdraɪv ə ˈkɑːr}{conduire une voiture}
\mot{speak English}{ˌspiːk ˈɪŋɡlɪʃ}{parler anglais}
\end{vocabulaire}

\begin{vocabulaire}{Dire ce qu'on sait faire}{12.06}
\categorie{Questions et réponses}
\mot{Can you swim?}{kən juː ˈswɪm}{Tu sais nager ?}
\mot{Yes, I can.}{ˌjes aɪ ˈkæn}{Oui. (réponse courte)}
\mot{No, I can't.}{ˌnoʊ aɪ ˈkænt}{Non. (réponse courte)}
\mot{I can sing.}{aɪ kən ˈsɪŋ}{Je sais chanter.}
\mot{I can't ski.}{aɪ ˈkænt ˈskiː}{Je ne sais pas skier.}
\categorie{Le niveau}
\mot{very well}{ˌveri ˈwel}{très bien}
\mot{quite well}{ˌkwaɪt ˈwel}{plutôt bien}
\mot{not very well}{ˌnɑːt ˌveri ˈwel}{pas très bien}
\mot{not at all}{ˌnɑːt ət ˈɑːl}{pas du tout}
\mot{fast}{fæst}{vite}
\categorie{Le sport}
\mot{a team}{ə ˈtiːm}{une équipe}
\mot{a game}{ə ˈɡeɪm}{un match, un jeu}
\mot{a band}{ə ˈbænd}{un groupe (de musique)}
\mot{sick}{sɪk}{malade}
\mot{Of course I can!}{əv ˌkɔːrs aɪ ˈkæn}{Bien sûr que je peux !}
\end{vocabulaire}

%% ---------------------------------------------------------------------
\rubrique{grammaire}{Can et les adverbes}

\otIntertitre[otViolet]{1 · \en{can} : la même forme pour tous}
\begin{center}\otTableau
\begin{tabular}{@{}l@{\hspace{8mm}}>{\enOn}l<{\enOff}@{\hspace{8mm}}l@{}}
\otTH{Forme} & \otTH{Exemple} & \otTH{Français}\\\midrule
affirmation      & I / you / he / she / we / they \textbf{can} swim. & je sais nager…\\
négation         & She \textbf{can't} (cannot) ski.                  & elle ne sait pas skier\\
question         & \textbf{Can} you play the piano?                  & tu sais jouer du piano ?\\
réponses courtes & Yes, I \textbf{can}. / No, I \textbf{can't}.      & oui / non\\\bottomrule
\end{tabular}\end{center}

\begin{regle}[Trois règles d'or]
\begin{etapes}
\item \en{can} ne prend \textbf{jamais} de \en{-s} : \en{she can}, pas \en{she cans}.
\item Il est suivi de la base verbale \textbf{sans} \en{to} : \en{I can swim}, pas \en{I can to swim}.
\item Pas de \en{do} ni de \en{does} : \en{Can you…?}, \en{He can't…}
\end{etapes}
\end{regle}

\blocaudio{12.07}{Écoutez et répétez}{Six phrases avec \en{can} et \en{can't}.}

\otIntertitre[otViolet]{2 · Dire son niveau}
\begin{center}\otTableau
\begin{tabular}{@{}c@{\hspace{8mm}}>{\enOn}l<{\enOff}@{\hspace{8mm}}l@{}}
\otEtoiles{3} & I can swim very well.            & très bien\\
\otEtoiles{2} & I can swim quite well.           & plutôt bien\\
\otEtoiles{1} & I can swim, but not very well.   & pas très bien\\
\otEtoiles{0} & I can't swim at all.             & pas du tout\\\bottomrule
\end{tabular}\end{center}

\begin{regle}[\en{good} ou \en{well} ?]
\en{good} est un adjectif (il décrit un nom) : \en{a good swimmer} (un bon nageur). \en{well} est un
adverbe (il décrit un verbe) : \en{She swims well} (elle nage bien). Ne dites jamais \en{She swims good}.
\end{regle}

\blocaudio{12.08}{Écoutez et répétez}{Une question avec \en{can}, puis sa réponse courte et une précision.}

\begin{attention}[Les pièges des francophones]
\fauxjuste{She cans swim.}{She can swim.}
\fauxjuste{I can to drive.}{I can drive.}
\fauxjuste{Do you can ski?}{Can you ski?}
\fauxjuste{I can swim very good.}{I can swim very well.}
\fauxjuste{He don't can sing.}{He can't sing.}
\end{attention}

%% ---------------------------------------------------------------------
\rubrique[can ou can't]{sons}{\en{can} ou \en{can't} ?}

En américain, la différence entre \en{can} et \en{can't} est subtile : le \en{t} final de \en{can't}
est souvent à peine prononcé. Fiez-vous à deux indices : dans une phrase affirmative, \en{can} est
\textbf{faible et rapide} \phon{kən}, alors que \en{can't} est \textbf{accentué}, avec un \phon{æ} bien
ouvert \phon{kænt}. Au Royaume-Uni, c'est plus simple : \en{can't} se dit \phon{kɑːnt}.

\begin{paires}{12.09}{Écoutez la phrase affirmative, puis la négative, et répétez en accentuant bien
\en{can't}.}
\sonpaire{I can swim.}{aɪ kən ˈswɪm}{je sais nager}{I \acc{can't} swim.}{aɪ ˈkænt ˈswɪm}{je ne sais pas nager}
\sonpaire{She can sing.}{ʃi kən ˈsɪŋ}{elle sait chanter}{She \acc{can't} sing.}{ʃi ˈkænt ˈsɪŋ}{elle ne sait pas chanter}
\sonpaire{We can come.}{wi kən ˈkʌm}{nous pouvons venir}{We \acc{can't} come.}{wi ˈkænt ˈkʌm}{nous ne pouvons pas venir}
\sonpaire{He can drive.}{hi kən ˈdraɪv}{il sait conduire}{He \acc{can't} drive.}{hi ˈkænt ˈdraɪv}{il ne sait pas conduire}
\sonpaire{You can dance.}{ju kən ˈdæns}{tu sais danser}{You \acc{can't} dance.}{ju ˈkænt ˈdæns}{tu ne sais pas danser}
\sonpaire{They can ski.}{ðeɪ kən ˈskiː}{ils savent skier}{They \acc{can't} ski.}{ðeɪ ˈkænt ˈskiː}{ils ne savent pas skier}
\end{paires}

%% ---------------------------------------------------------------------
\rubrique{culture}{Le sport aux États-Unis}

\begin{culture}[Les quatre grands sports]
Le \en{football} américain (on joue surtout avec les mains !), le \en{baseball}, surnommé
\en{America's pastime}, le \en{basketball} et le hockey sur glace sont les sports rois. Notre
football s'appelle \en{soccer} aux États-Unis. La finale du football américain, le \en{Super Bowl},
est l'émission la plus regardée de l'année.
\end{culture}

\begin{paire}
\begin{culture}[New York et ses équipes]
Baseball : les \en{Yankees} et les \en{Mets} ; basket : les \en{Knicks} et les \en{Brooklyn Nets} ;
hockey : les \en{Rangers}. Les \en{Giants} et les \en{Jets}, équipes de football américain, jouent… dans
le New Jersey voisin !
\end{culture}
\begin{culture}[Le sport à l'école]
Aux États-Unis, le sport universitaire (\en{college sports}) est très populaire : les matchs remplissent
des stades immenses, et les meilleurs étudiants-athlètes obtiennent des bourses d'études.
\end{culture}
\end{paire}

%% ---------------------------------------------------------------------
\rubrique{pratique}{Exercices}

\exercice[ecrire,niveau=1]{Les talents}{Écrivez l'expression anglaise.}
\begin{questions}[2][french]
\item nager → \enb{\trou[34mm]{swim}}
\item faire du vélo → \enb{\trou[34mm]{ride a bike}}
\item jouer du piano → \enb{\trou[34mm]{play the piano}}
\item conduire une voiture → \enb{\trou[34mm]{drive a car}}
\item chanter → \enb{\trou[34mm]{sing}}
\item jouer au tennis → \enb{\trou[34mm]{play tennis}}
\item dessiner → \enb{\trou[34mm]{draw}}
\item courir → \enb{\trou[34mm]{run}}
\end{questions}

\exercice[lire,ecrire,niveau=1]{Qui sait faire quoi ?}{Regardez le tableau et complétez avec \en{can}
ou \en{can't}.}
\otTalents
\begin{questions}[2]
\item Léa \trou[16mm]{can} swim.
\item Jack \trou[16mm]{can't} ski.
\item Emma \trou[16mm]{can't} swim.
\item Jack \trou[16mm]{can} play the guitar.
\item Léa \trou[16mm]{can't} play the guitar.
\item Emma \trou[16mm]{can} cook.
\item Jack \trou[16mm]{can't} cook.
\item Emma \trou[16mm]{can} ski.
\end{questions}

\exercice[lire,ecrire,niveau=2]{Questions et réponses}{Toujours d'après le tableau, répondez par une
réponse courte.}
\begin{questions}
\item Can Léa cook? \hfill \trou[40mm]{Yes, she can.}
\item Can Jack ski? \hfill \trou[40mm]{No, he can't.}
\item Can Emma play the guitar? \hfill \trou[40mm]{No, she can't.}
\item Can Jack swim? \hfill \trou[40mm]{Yes, he can.}
\item Can Léa and Emma ski? \hfill \trou[40mm]{Yes, they can.}
\end{questions}

\exercice[ecrire,niveau=2]{Quel niveau ?}{Complétez avec \en{very well}, \en{quite well},
\en{not very well}, \en{at all} ou \en{very fast}, d'après les étoiles.}
\begin{questions}
\item \otEtoiles{3}\quad Sam plays the guitar \trou[30mm]{very well}.
\item \otEtoiles{2}\quad Léa plays the piano \trou[30mm]{quite well}.
\item \otEtoiles{1}\quad Léa can sing, but \trou[30mm]{not very well}.
\item \otEtoiles{0}\quad Jack can't ski \trou[30mm]{at all}.
\item \otEtoiles{3}\quad Emma runs \trou[30mm]{very fast}.
\end{questions}

\exercice[ecrire,niveau=2]{Corrigez l'erreur}{Chaque phrase contient une erreur. Récrivez-la correctement.}
\begin{questions}
\item She cans swim. \hfill \trou[66mm]{She can swim.}
\item I can to drive. \hfill \trou[66mm]{I can drive.}
\item Do you can ski? \hfill \trou[66mm]{Can you ski?}
\item He sings very good. \hfill \trou[66mm]{He sings very well.}
\item He don't can dance. \hfill \trou[66mm]{He can't dance.}
\end{questions}

\exercice[lire,niveau=2]{L'annonce de Sam}{Lisez l'annonce, puis indiquez si chaque phrase est vraie
(V) ou fausse (F).}
\begin{anglais}
\begin{tcolorbox}[otBase,colback=white,colframe=otRule,boxrule=0.5pt,fontupper=\fontsize{10.8}{14.4}\selectfont]
{\otSemi\fontsize{11}{13}\selectfont SINGER WANTED!}\par\vspace{1.4mm}
We're a jazz band from Brooklyn: Sam (guitar), Ben (drums) and Kim (piano). Our singer is sick!
Can you sing jazz? Can you come on Friday at 8 p.m.? You don't need to read music, but you need to
sing well and speak English. Call Sam: 718-555-0162.
\par\vspace{1mm}{\fontsize{8.6}{10}\selectfont\color{otMuted}\foreignlanguage{french}{\textit{wanted} : recherché · \textit{drums} : batterie · \textit{read music} : lire la musique · \textit{need} : avoir besoin}}
\end{tcolorbox}
\begin{questions}[2]
\item Sam plays the guitar. \hfill\trou[9mm]{V}
\item Kim plays the drums. \hfill\trou[9mm]{F}
\item They need a singer. \hfill\trou[9mm]{V}
\item The concert is on Saturday. \hfill\trou[9mm]{F}
\item You need to read music. \hfill\trou[9mm]{F}
\item You need to sing well. \hfill\trou[9mm]{V}
\end{questions}
\end{anglais}

\needspace{12\baselineskip}% en bas de page, l'exercice débordait de 1,4 pt : il passe à la page suivante
\exercice[ecouter,niveau=1,piste=12.10]{\en{can} ou \en{can't} ?}{Écoutez chaque phrase, dite deux
fois. Cochez ce que vous entendez.}
\begin{questions}[2]
\item Léa \hfill \choix{can}\quad\choix{can't}\quad ski.
\item Jack \hfill \choix{can}\quad\choix{can't}\quad cook.
\item Emma \hfill \choix{can}\quad\choix{can't}\quad swim.
\item Sam \hfill \choix{can}\quad\choix{can't}\quad sing.
\item Maya \hfill \choix{can}\quad\choix{can't}\quad drive.
\item Mr.\ Carter \hfill \choix{can}\quad\choix{can't}\quad dance.
\end{questions}

\exercice[ecouter,ecrire,niveau=2,piste=12.11]{Les talents de Maya}{Sam interroge sa sœur Maya.
Écoutez et écrivez \en{yes} ou \en{no}.}
\begin{center}\otTableau
\begin{tabular}{@{}l@{\hspace{16mm}}c@{}}
\otTH{Maya can\ldots} & \otTH{yes / no}\\\midrule
\enb{sing}            & \enb{\trou[16mm]{yes}}\\[0.6mm]
\enb{play the piano}  & \enb{\trou[16mm]{no}}\\[0.6mm]
\enb{dance}           & \enb{\trou[16mm]{yes}}\\[0.6mm]
\enb{drive a car}     & \enb{\trou[16mm]{no}}\\[0.6mm]
\enb{speak Spanish}   & \enb{\trou[16mm]{yes}}\\\bottomrule
\end{tabular}\end{center}

\exercice[ecouter,ecrire,niveau=2,piste=12.12]{Dictée à trous}{Le coach parle avec Emma. Écoutez et
complétez.}
\begin{anglais}
\noindent\begin{tabular}{@{}l@{\hspace{6mm}}l@{}}
\otNomDit{coach} & \trou{Can} you play tennis, Emma?\\[1.8mm]
\otNomDit{emma}  & Yes, I \trou{can}. I play quite \trou{well}.\\[1.8mm]
\otNomDit{coach} & And can you \trou{swim}?\\[1.8mm]
\otNomDit{emma}  & No, I \trou{can't}. I can't swim at \trou{all}!\\[1.8mm]
\otNomDit{coach} & No problem. Can you \trou{run}?\\[1.8mm]
\otNomDit{emma}  & Oh yes! I run very \trou{fast}!\\
\end{tabular}
\end{anglais}

\exercice[ecrire,niveau=2]{Traduisez}{Traduisez en anglais.}
\begin{questions}[1][french]
\item Je sais nager, mais je ne sais pas skier.\lignes{1}
\item Tu sais jouer de la guitare ? — Oui. \textit{(réponse courte)}\lignes{1}
\item Elle chante très bien.\lignes{1}
\item Il ne sait pas conduire.\lignes{1}
\item Tu peux venir vendredi ?\lignes{1}
\item Je joue du piano, mais pas très bien.\lignes{1}
\item Ils ne savent pas danser du tout.\lignes{1}
\item C'est une bonne nageuse.\lignes{1}
\end{questions}

\exercice[ecrire,niveau=3]{Vos talents}{En six phrases, dites ce que vous savez faire (et à quel
niveau), et ce que vous ne savez pas faire.}
\lignes{6}

\exercice[ecouter,parler,niveau=3,piste=12.13]{À vous de parler !}{Le coach vous interroge à votre tour.
Répondez \textbf{à voix haute} pendant chaque pause.}
\begin{center}\otTableau
\begin{tabular}{@{}>{\otSemi\fontsize{8.6}{10}\selectfont\color{otNavy}}l@{\hspace{3mm}}>{\enOn}p{70mm}<{\enOff}@{\hspace{6mm}}p{78mm}@{}}
 & \otTH{Le coach dit\ldots} & \otTH{Vous\ldots}\\\midrule
1. & Can you swim? & répondez, et précisez votre niveau.\\
2. & Can you play basketball? & répondez par une réponse courte.\\
3. & What sports can you do? & citez un ou deux sports.\\
4. & Can you play a musical instrument? & répondez et dites lequel.\\
5. & Great! Can you come next Saturday? & dites que oui.\\\bottomrule
\end{tabular}\end{center}

%% ---------------------------------------------------------------------
\rubrique{bilan}{Faites le point}

\begin{essentiel}
\begin{minipage}[t]{0.54\linewidth}
\otEssTitre{Les phrases clés}
\renewcommand{\arraystretch}{1.32}
\begin{tabular}{@{}>{\enOn\otSemi\fontsize{9.4}{11}\selectfont}l<{\enOff}@{\hspace{3mm}}>{\fontsize{9.2}{11}\selectfont\color{white!78!otNavy}}l@{}}
Can you swim?              & Tu sais nager ?\\
Yes, I can. / No, I can't. & Oui. / Non.\\
I can sing very well.      & Je chante très bien.\\
I can't ski at all.        & Je ne sais pas du tout skier.\\
Can you come on Friday?    & Tu peux venir vendredi ?\\
She's a good swimmer.      & C'est une bonne nageuse.\\
\end{tabular}
\end{minipage}\hfill
\begin{minipage}[t]{0.41\linewidth}
\otEssTitre{\en{can}}
{\fontsize{9.2}{11.6}\selectfont même forme pour tous · base verbale sans \en{to} · pas de \en{do}\par}
\otEssTitre{Le niveau}
{\enb{\fontsize{9.2}{11.6}\selectfont very well · quite well · not very well · not at all}\par
{\fontsize{9}{11.4}\selectfont\color{white!78!otNavy}\en{good} (adjectif) ≠ \en{well} (adverbe)}\par}
\otEssTitre{Prononciation}
{\fontsize{9.2}{11.6}\selectfont \en{I can swim} \phon{kən} · \en{I can't swim} \phon{ˈkænt}\par}
\end{minipage}
\end{essentiel}

\begin{jesais}
\item dire ce que je sais faire et ce que je ne sais pas faire ;
\item poser une question avec \en{can} et répondre par une réponse courte ;
\item préciser mon niveau avec \en{very well}, \en{quite well}, \en{not very well}, \en{not at all} ;
\item employer \en{well} (adverbe) et \en{good} (adjectif) à bon escient ;
\item nommer une quinzaine de sports et de loisirs ;
\item distinguer \en{can} et \en{can't} à l'oral.
\end{jesais}

\clearpage
\begin{testexpress}
\qcm[b]{« Elle sait nager » se dit :}{\en{She cans swim.}}{\en{She can swim.}}{\en{She can to swim.}}
\qcm[c]{« Tu sais conduire ? » se dit :}{\en{Do you can drive?}}{\en{You can drive?}}{\en{Can you drive?}}
\qcm[a]{\en{Can you ski?} — Réponse négative :}{\en{No, I can't.}}{\en{No, I don't.}}{\en{No, I can.}}
\qcm[b]{« Il joue très bien » se dit :}{\en{He plays very good.}}{\en{He plays very well.}}{\en{He very well plays.}}
\qcm[a]{« jouer du piano » se dit :}{\en{play the piano}}{\en{play piano of}}{\en{play at the piano}}
\qcm[c]{« jouer au tennis » se dit :}{\en{play the tennis}}{\en{play at tennis}}{\en{play tennis}}
\qcm[b]{« pas du tout » se dit :}{\en{not very}}{\en{not at all}}{\en{not all}}
\qcm[a]{Dans \en{I can swim}, \en{can} se prononce :}{\phon{kən}, faible}{\phon{kænt}}{\phon{kɑːn}}
\qcm[c]{Le \en{soccer}, c'est :}{le football américain}{le hockey}{notre football}
\qcm[b]{« une bonne nageuse » se dit :}{\en{a well swimmer}}{\en{a good swimmer}}{\en{a swimmer good}}
\end{testexpress}

\noindent{\fontsize{9.6}{12}\selectfont\color{otMuted}Niveau visé (CECRL, niveau A1) : « peut dire ce
qu'il sait faire et poser des questions simples sur les capacités des autres ». Les corrigés de
l'unité sont page~\pageref{cor-u12-1} ; les textes des pistes d'exercices,
page~\pageref{trans-12.10}.\par}
